
\medskip
\phantomsection\label{prop:sga3-VIB-8.4}%
\noindent\textbf{Proposition 8.4.---}
\emph{Let \(G\) be an \(S\)-group of finite presentation such
that, for every \(s\in S\), the fiber \(G_s\) is smooth over
\(\kappa(s)\). Let \(T\) be the set of points \(s\in S\) for
which \(G_s\) is solvable \textup{(respectively, nilpotent)}.
Then \(T\) is locally constructible in \(S\). If, in addition,
\(G\) is assumed flat and separated over \(S\) \textup{(that is,
when \(G\) is smooth, of finite presentation and separated over
\(S\))}, then \(T\) is closed in \(S\).}

\emph{Proof.}
We may clearly suppose that \(S\) is affine, with ring \(A\).
By \hyperref[par:sga3-VIB-10.1]{10.1} and
\hyperref[par:sga3-VIB-10.10]{10.10\textup{b)}},%
\begingroup
\renewcommand{\thefootnote}{\fnsymbol{footnote}}%
\footnote[1]{\label{note:sga3-VIB-8.4-star}
In the course of this proof we shall use results established
in \hyperref[sec:sga3-VIB-10]{Section~10}; like \hyperref[sec:sga3-VIB-9]{Section~9}, they do not depend on the
present \hyperref[sec:sga3-VIB-8]{Section~8}.}%
\endgroup\ 
there exist a Noetherian subring \(A'\) of \(A\) and an
\(A'\)-group \(G'\) of finite type such that
\[
  G'\otimes_{A'}A\simeq G.
\]
Moreover (EGA~IV\(_3\), 8.10.5 and 11.2.6),\footnote{Translator's note: the 1970 text cites ``EGA IV
10.8.5''; the reference is EGA~IV\(_3\), 8.10.5~(v) (separation in a projective limit), as corrected in the
Gille--Polo re-edition (2011), Exposé~VI\textsubscript{B}.}
if \(G\) is flat and separated over \(S\), we may suppose that
\(G'\) is flat and separated over
\(S'=\operatorname{Spec}A'\).%

Likewise (EGA~IV\(_3\), 9.7.7 and 9.3.3), we may suppose that,
for every \(s'\in S'\), \(G'_{s'}\) is smooth over
\(\kappa(s')\).\footnote{Translator's note: as the Gille--Polo re-edition (2011), Exposé~VI\textsubscript{B},
explains, the fibers of a group of finite presentation are smooth exactly where they are geometrically
reduced (see the note to \hyperref[concl:sga3-VIB-7.1-i]{7.1~(i)}), and the set of such points is locally
constructible (EGA~IV\(_3\), 9.7.7); as it contains every point of \(S\), it contains every point of
\(\operatorname{Spec}A'_\lambda\) for a suitable finitely generated subring \(A'_\lambda\supset A'\) of \(A\),
by EGA~IV\(_3\), 9.3.3, and one replaces \(A'\) by \(A'_\lambda\).} On the other hand,
if \(s'\) denotes the image of \(s\) in \(S'\), then
\[
  G'_{s'}\otimes_{\kappa(s')}\kappa(s)\simeq G_s.
\]
It follows from \hyperref[prop:sga3-VIB-8.3]{8.3} that \(G_s\)
is solvable (respectively, nilpotent) if and only if the same
is true of \(G'_{s'}\). We are therefore reduced to the case
where \(S\) is a Noetherian affine scheme.

We first show that \(T\) is constructible. Applying the
criterion of EGA~0\(_{\mathrm{III}}\), 9.2.3, and arguing as
above, we are reduced to showing that, when \(S\) is
Noetherian and integral, either \(T\) or \(S-T\) contains a
nonempty open subset of \(S\).

Suppose, then, that \(S\) is integral and Noetherian, with
generic point \(\eta\). For every \(s\in S\), put
\(K_s^0=G_s\), and put, for \(p\geq1\),
\phantomsection\label{form:sga3-VIB-8.4-Ksp}%
\[
  K_s^p=(K_s^{p-1},K_s^{p-1})
  \qquad
  \bigl(\textup{respectively, }\
    K_s^p=(G_s,K_s^{p-1})\bigr).
\tag*{\(\dagger\)}
\]
\begingroup
\renewcommand{\thefootnote}{\fnsymbol{footnote}}%
\footnotetext[2]{\label{note:sga3-VIB-8.4-G-vs-Gs}
Translator's note: in the nilpotent branch of this display,
the French source prints \(G\). Since \(K_s^{p-1}\) is a
\(\kappa(s)\)-group and the subsequent formulas use \(G_s\),
the fiber subscript is restored here.}%
\endgroup
We first show that the sequence of closed subschemes
\(K_\eta^p\) is stationary. By
\hyperref[concl:sga3-VIB-7.3-v]{7.3\textup{(v)}}, each
\(K_\eta^p\) is invariant, and hence the sequence
\((K_\eta^p)_p\) is decreasing. It is therefore stationary,
since \(G_\eta\) is Noetherian. Thus there is an integer \(n\)
such that
\phantomsection\label{form:sga3-VIB-8.4-stationary}%
\[
  K_\eta^p=K_\eta^n
  \qquad\text{for every }p\geq n.
\]

On the other hand, we shall see in
\hyperref[prop:sga3-VIB-10.12]{10.12} and
\hyperref[cor:sga3-VIB-10.13]{10.13} that there exist a
nonempty open subset \(S'\) of \(S\) and an \(S'\)-group \(D\)
of finite presentation such that, for every \(s\in S'\),
\phantomsection\label{form:sga3-VIB-8.4-D-fibers}%
\[
  D_s=K_s^n,
  \qquad
  (D_s,D_s)=D_s
  \quad
  \bigl(\textup{respectively, }(G_s,D_s)=D_s\bigr).
\]
We may suppose that \(S'=S\). Then, for every \(s\in S\) and
every \(p\geq n\), one has \(D_s=K_s^p\). Consequently, \(G_s\)
is solvable (respectively, nilpotent) if and only if \(D_s\)
is isomorphic to the unit \(\kappa(s)\)-group.

But by EGA~IV\(_3\), 9.6.1\textup{(xi)}, the set of points
\(s\in S\) for which the structural morphism
\[
  D_s\longrightarrow\operatorname{Spec}\kappa(s)
\]
is an isomorphism is constructible,\footnote{Translator's note: EGA~IV\(_3\), 9.6.1 gives a locally
constructible set; it is constructible because \(S\) is affine, hence quasi-compact with a basis of
retrocompact open subsets (EGA~\(0_{\mathrm{III}}\), 9.1.12), as the Gille--Polo re-edition (2011),
Exposé~VI\textsubscript{B}, notes.} and it is therefore either
nowhere dense or contains a nonempty open subset of \(S\).

We now show that, if \(G\) is flat and separated over \(S\),
then \(T\) is closed. Since \(T\) is locally
constructible, it is closed if and only if it is stable under
specialization (EGA~IV\(_1\), 1.10.1).

Let \(s\in S\), and let \(s'\) be a specialization of \(s\)
in \(S\). Since we have reduced to the case where \(S\) is
Noetherian, EGA~II, 7.1.9 provides a discrete valuation ring
\(A\) and a morphism
\[
  \operatorname{Spec}A\longrightarrow S
\]
such that \(s\), respectively \(s'\), is the image of the
generic point \(\alpha\), respectively the closed point \(a\),
of \(\operatorname{Spec}A\). It is enough to show that, if
\[
  G'=G\times_S\operatorname{Spec}A
\]
and \(G'_\alpha\) is solvable (respectively, nilpotent), then
the same is true of \(G'_a\). Since \(G\) is flat and
separated over \(S\), the \(A\)-group \(G'\) is flat and
separated.%
\begingroup
\renewcommand{\thefootnote}{\fnsymbol{footnote}}%
\footnote[3]{\label{note:sga3-VIB-8.4-Gprime-alpha}
Translator's note: the French source prints \(G'_\alpha\) as
being flat and separated over \(A\). The intended object is
the base-changed \(A\)-group \(G'\); its generic fiber
\(G'_\alpha\) is a group over \(\kappa(\alpha)\).}%
\endgroup\ 
We are thus reduced to the case where \(S\) is the spectrum of
a discrete valuation ring \(A\).

Since \(G_\alpha\) is assumed solvable (respectively,
nilpotent), there is an integer \(q\) such that
\(K_\alpha^q\), with the notation introduced above, is
isomorphic to the unit \(\kappa(\alpha)\)-group. For every
\(n\), let \(\overline{K_\alpha^n}\) denote the schematic
closure of \(K_\alpha^n\) in \(G\), in the sense of
EGA~IV\(_2\), 2.8.5. We show by induction on \(p\) that
\phantomsection\label{form:sga3-VIB-8.4-closure-induction}%
\[
  \bigl(\overline{K_\alpha^p}\bigr)_a
    \supset K_a^p.
\]
First, since \(G\) is flat over \(A\), one has
\(\overline{G_\alpha}=G\) by EGA~IV\(_2\), 2.8.5, and hence
\[
  \bigl(\overline{K_\alpha^0}\bigr)_a=K_a^0.
\]

Denote by \(\overline{\nu},\overline{\nu}_a,\nu_a,\nu_\alpha\)
the morphisms
\phantomsection\label{form:sga3-VIB-8.4-nu-solvable}%
\[
\begin{aligned}
  \overline{\nu}:\;&
    \overline{K_\alpha^p}\times_A\overline{K_\alpha^p}
    \longrightarrow G,\\
  \overline{\nu}_a:\;&
    \bigl(\overline{K_\alpha^p}\bigr)_a
    \times_{\kappa(a)}
    \bigl(\overline{K_\alpha^p}\bigr)_a
    \longrightarrow G_a,\\
  \nu_a:\;&K_a^p\times_{\kappa(a)}K_a^p
    \longrightarrow G_a,\\
  \nu_\alpha:\;&K_\alpha^p\times_{\kappa(\alpha)}K_\alpha^p
    \longrightarrow G_\alpha
\end{aligned}
\]
respectively
\phantomsection\label{form:sga3-VIB-8.4-nu-nilpotent}%
\[
\begin{aligned}
  \overline{\nu}:\;&G\times_A\overline{K_\alpha^p}
    \longrightarrow G,\\
  \overline{\nu}_a:\;&
    G_a\times_{\kappa(a)}
    \bigl(\overline{K_\alpha^p}\bigr)_a
    \longrightarrow G_a,\\
  \nu_a:\;&G_a\times_{\kappa(a)}K_a^p
    \longrightarrow G_a,\\
  \nu_\alpha:\;&G_\alpha\times_{\kappa(\alpha)}K_\alpha^p
    \longrightarrow G_\alpha,
\end{aligned}
\]
defined as in \hyperref[def:sga3-VIB-7.2-vii]{7.2\textup{(vii)}}. Then, since \(\nu_\alpha\) factors through \(K_\alpha^{p+1}\),
\(\overline{\nu}\) factors through
\(\overline{K_\alpha^{p+1}}\). The latter is clearly a group
subscheme of \(G\), and therefore contains
\(\Gamma_G(\overline{\nu})\). By
\hyperref[concl:sga3-VIB-7.1-iii]{7.1\textup{(iii)}},
\[
  \Gamma_G(\overline{\nu})_a
    =\Gamma_{G_a}(\overline{\nu}_a).
\]
Moreover, the induction hypothesis gives
\phantomsection\label{form:sga3-VIB-8.4-induction-inclusions}%
\[
\begin{aligned}
  K_a^p\times_{\kappa(a)}K_a^p
  &\subset
  \bigl(\overline{K_\alpha^p}\bigr)_a
  \times_{\kappa(a)}
  \bigl(\overline{K_\alpha^p}\bigr)_a
  &&\text{in the solvable case},\\
  G_a\times_{\kappa(a)}K_a^p
  &\subset
  G_a\times_{\kappa(a)}
  \bigl(\overline{K_\alpha^p}\bigr)_a
  &&\text{in the nilpotent case}.
\end{aligned}
\]
Consequently,
\phantomsection\label{form:sga3-VIB-8.4-Gamma-chain}%
\[
  K_a^{p+1}
  =\Gamma_{G_a}(\nu_a)
  \subset\Gamma_{G_a}(\overline{\nu}_a)
  \subset\bigl(\overline{K_\alpha^{p+1}}\bigr)_a.
\]

Finally, \(K_\alpha^q\) is isomorphic to the unit
\(\kappa(\alpha)\)-group. The unit \(A\)-group is flat over
\(A\) and is isomorphic to a closed subscheme of \(G\), since
\(G\) is separated over \(A\)
\textup{(cf.\ \hyperref[prop:sga3-VIB-5.1]{5.1})}. It follows
from EGA~IV\(_2\), 2.8.5 that
\(\overline{K_\alpha^q}\) is isomorphic to the unit
\(A\)-group. Since we have proved that
\[
  K_a^q\subset\bigl(\overline{K_\alpha^q}\bigr)_a,
\]
the group \(K_a^q\) is isomorphic to the unit
\(\kappa(a)\)-group. Thus \(G_a\) is solvable
\textup{(respectively, nilpotent)}.
